\documentclass[12pt]{article}

\usepackage{amsmath}
\usepackage{amssymb}
\usepackage{amsthm}
\usepackage{etoolbox}
\usepackage{hyperref}
\usepackage{geometry}
\usepackage{draftwatermark}
\SetWatermarkText{DRAFT}
\SetWatermarkScale{1.5}
\SetWatermarkColor[gray]{0.85}

\geometry{margin=1in}

%------------------------------------------------------------
% Theorem environments
%------------------------------------------------------------
\newtheorem{theorem}{Theorem}[section]
\newtheorem{lemma}[theorem]{Lemma}
\newtheorem{proposition}[theorem]{Proposition}
\newtheorem{corollary}[theorem]{Corollary}
\theoremstyle{definition}
\newtheorem{definition}[theorem]{Definition}
\newtheorem{example}[theorem]{Example}
\theoremstyle{remark}
\newtheorem{remark}[theorem]{Remark}
\theoremstyle{plain}
\newtheorem{conjecture}[theorem]{Conjecture}

%------------------------------------------------------------
% Notation macros
%------------------------------------------------------------
\newcommand{\ZZ}{\mathbb{Z}}
\newcommand{\NN}{\mathbb{N}}
\newcommand{\RR}{\mathbb{R}}
\newcommand{\OO}{\mathsf{O}}
\newcommand{\EE}{\mathsf{E}}

%------------------------------------------------------------
% Hyperref setup
%------------------------------------------------------------
\hypersetup{
  colorlinks=true,
  linkcolor=blue,
  citecolor=blue,
  urlcolor=blue,
}

\title{Mediants and Collatz:\\
  Constructing Generalised Collatz Trajectories\\
  via Rational Approximation}
\author{Jon Seymour}
\date{2026}

\begin{document}

\maketitle

\begin{abstract}
We describe an application of the mediant tree algorithm of~\cite{paper76} to the
construction of generalised Collatz trajectories.
Given a generalised Collatz system $(g, h)$ with odd step $x \mapsto gx + q$ and
even step $x \mapsto x/h$, a parity sequence realising a good rational approximant
to $\theta = \log_h g$ is produced directly by the mediant tree as a mediant word
over $\{\OO, \EE\}$.
We show how to convert such a mediant word into a concrete starting integer $x_0$
via a standard modular-arithmetic pipeline, and explain why the canonical greedy
parity-sequence algorithm produces exactly the Sturmian word for slope
$1/(1+\log_h g)$.
We compare our approach to Eliahou's continued-fraction bounding
method~\cite{eliahou1993}, and note that the mediant tree simultaneously produces
both the rational approximant and a concrete trajectory realising it --- a capability
absent from Eliahou's framework.
\end{abstract}

\tableofcontents

%------------------------------------------------------------
\section{Introduction}
\label{sec:intro}
%------------------------------------------------------------

A \emph{generalised Collatz sequence} is defined by a pair $(g, h)$ with $g, h \geq 2$
and $\gcd(g, h) = 1$: an odd element $x$ is mapped to $gx + q$ and an even element
$x$ is divided by $h$.
The standard Collatz sequence corresponds to $g = 3$, $h = 2$, $q = 1$.
The parity structure of a trajectory --- the sequence of odd and even steps --- is
governed by $\theta = \log_h g$: a long-descending trajectory must have the ratio
of even to odd steps approximate $\theta$ from above.

The companion paper~\cite{paper76} develops the \emph{mediant tree}, an algorithm
that constructs increasingly precise rational approximants to any real $\theta > 1$
together with a binary string over $\{\OO, \EE\}$ --- the \emph{mediant word} ---
that encodes the approximation path.
The mediant word is proven to be identical to the Sturmian word with slope
$1/(1+\theta)$ at all interior positions, establishing a tight connection between
rational approximation and combinatorics on words.

The present paper describes the Collatz application that originally motivated
the development of~\cite{paper76}.
The main contributions are:
\begin{enumerate}
  \item A pipeline (Section~\ref{sec:pipeline}) for converting a mediant word into
    a concrete starting integer $x_0$ whose generalised Collatz trajectory follows
    the prescribed parity pattern for exactly $e + o$ steps.
  \item A description (Section~\ref{sec:greedy}) of the greedy parity-sequence
    algorithm and a proof that it produces the Sturmian word for slope
    $1/(1+\log_h g)$.
  \item A discussion (Section~\ref{sec:eliahou}) of the relationship to Eliahou's
    work on rational approximation of $\log_2 3$ in the context of Collatz cycle
    bounds.
\end{enumerate}

\paragraph{Organisation.}
Section~\ref{sec:prelim} recalls the key definitions and results from~\cite{paper76}.
Section~\ref{sec:collatz} introduces generalised Collatz sequences and parity
sequences.
Section~\ref{sec:pipeline} presents the step-by-step pipeline from mediant word
to Collatz trajectory.
Section~\ref{sec:greedy} describes the greedy algorithm and its Sturmian
interpretation.
Section~\ref{sec:not-sturmian} discusses why general Collatz trajectories are not
Sturmian.
Section~\ref{sec:eliahou} compares with Eliahou's approach.

%------------------------------------------------------------
\section{Preliminaries}
\label{sec:prelim}
%------------------------------------------------------------

We recall the definitions and results from~\cite{paper76} that are used in this
paper.
All notation is consistent with that reference.

A \emph{ratio} is a pair $(e, o) \in \NN^2$ with $\gcd(e, o) = 1$, written $e/o$.
Given a real $\theta > 1$, a ratio $r = e/o$ is a \emph{good estimate} of $\theta$
if $|e - \theta o| < 1$.
Two good estimates $r_i = e_i/o_i$ and $r_{i+1} = e_{i+1}/o_{i+1}$ form a
\emph{bounded pair} $(r_i, r_{i+1})$ if they bracket $\theta$ and satisfy
$\Delta(r_i, r_{i+1}) := e_{i+1} o_i - e_i o_{i+1} = \pm 1$.

Given a bounded pair $(r_i, r_{i+1})$ and a weight $w \geq 1$, the
\emph{weighted mediant} is
\[
  \mathrm{Mediant}(r_i, r_{i+1}, w) = \frac{e_i + w\,e_{i+1}}{o_i + w\,o_{i+1}}.
\]
The \emph{optimal weight} $w^*$ is the unique integer such that the pair
$\bigl(\mathrm{Mediant}(r_i, r_{i+1}, w^*-1),\, \mathrm{Mediant}(r_i, r_{i+1}, w^*)\bigr)$
brackets $\theta$; it is given in closed form by Theorem~8.1 of~\cite{paper76}.

The \emph{mediant word} $\pi(r)$ of a ratio $r$ produced by the mediant tree is a
string over $\{\OO, \EE\}$ of length $e + o$ encoding the sequence of weighted-mediant
choices made to reach $r$ from the initial bounded pair.
It is a lower Christoffel word (Lemma~9.3 of~\cite{paper76}), begins with $\EE\EE$
for every non-leaf estimate (Theorem~12.4 of~\cite{paper76}), and is identical to
the Sturmian word for slope $1/(1+\theta)$ at all interior positions
(Theorem~12.6 of~\cite{paper76}).

%------------------------------------------------------------
\section{Generalised Collatz Sequences and Parity Sequences}
\label{sec:collatz}
%------------------------------------------------------------

A \emph{generalised Collatz sequence} uses an odd step $x \mapsto gx + q$ and an
even step $x \mapsto x/h$, where $g, h \in \NN$ with $g, h \geq 2$ and
$\gcd(g, h) = 1$, and $q \in \ZZ$ with $q \neq 0$.
The parity sequence structure depends only on $g$ and $h$ through
$\theta = \log_h g$; the value of $q$ affects only the concrete starting
residue $x_0$ (Step~4 of the pipeline in Section~\ref{sec:pipeline}).
A \emph{parity sequence} is a string over $\{\OO, \EE\}$ recording the sequence
of odd and even steps along a trajectory; $\OO$ denotes an odd step and $\EE$ an
even step.

For a trajectory to descend on average, the ratio of even steps to odd steps must
exceed $\log_h g$: if a segment of length $e + o$ has $e$ even steps and $o$ odd
steps, the net effect is multiplication by $g^o / h^e$, which is less than $1$ iff
$e/o > \log_h g$.
The boundary case $e/o = \lceil o\log_h g \rceil / o$ is a good estimate of
$\log_h g$, and corresponds to the upper-bound condition
\[
  h^{e-1} < g^o < h^e.
\]
A parity sequence realising a good estimate $e/o$ must therefore have exactly $o$
copies of $\OO$ and $e$ copies of $\EE$, with no two consecutive $\OO$'s.
$\OO\OO$-avoidance holds on three independent grounds: it is a property of Sturmian
sequences with $\alpha < 1/2$ (equivalently $\theta > 1$); it holds for $h = 2$
trajectories when $q$ is odd and $\gcd(g, h) = 1$; and it is maintained by the
mediant word generator by construction.
By contrast, for $h \neq 2$ arbitrary trajectories need not be $\OO\OO$-free.
Every $\OO\OO$-free parity sequence can be realised as a $gx + q$, $x/h$
trajectory for some $q$ determined by the sequence~\cite{paper71}.

%------------------------------------------------------------
\section{From Mediant Word to Collatz Trajectory: the Pipeline}
\label{sec:pipeline}
%------------------------------------------------------------

The mediant tree applies to any real $\theta > 1$, not only to logarithms.
The following pipeline converts a mediant word into a concrete Collatz starting
integer for any $(g, h)$ with $\theta = \log_h g$.

\medskip
\noindent\textbf{Step 1: Choose a real $\theta$ and an initial bounded pair.}
Select any real $\theta > 1$ and a Farey-adjacent bounded pair $(r_0, r_1)$ with
$r_0 < \theta < r_1$.
For the standard Collatz system take $\theta = \log_2 3$; for the generalised
system $(g, h)$ take $\theta = \log_h g$.
The construction applies equally to any irrational $\theta$ (e.g.\ $\pi$,
$\sqrt{2}$), in which case it produces rational approximants without a Collatz
interpretation.

\medskip
\noindent\textbf{Step 2: Run the mediant tree.}
Apply the optimal weight formula (Theorem~8.1 of~\cite{paper76}) to produce a
sequence of mediant words $\pi(r_{2n+1})$ of increasing length.
At level $n$, the word $\pi(r_{2n+1})$ has $o_{2n+1}$ odd steps and $e_{2n+1}$
even steps, and $e_{2n+1}/o_{2n+1}$ is a good estimate of $\theta$ satisfying
$h^{e-1} < g^o < h^e$ (upper bound) or $h^e < g^o < h^{e+1}$ (lower bound).

\medskip
\noindent\textbf{Step 3: Derive the path identity.}
For a chosen mediant word $\sigma = \pi(r_{2n+1})$, the function
\texttt{path\_identity($\sigma$)} constructs the \emph{basic path identity}
\[
  h^e\,x_n = g^o\,x_0 + q\,K(g, h),
\]
where $x_0$ is the starting value, $x_n$ is the value after $n = e + o$ steps
following $\sigma$, and $K(g,h)$ is the \emph{carry polynomial} --- a polynomial
in $g$ and $h$ whose coefficients record the cumulative effect of the additive
terms along the trajectory, with $q$ scaling the carry to the chosen additive
constant.
This identity holds for \emph{every} $x_0$; it is a universal statement about the
linear map defined by $\sigma$.
Adding $g^o h^{e+1} t$ to both sides for any integer $t$ gives the equivalent
\emph{shifted path identity}
\[
  h^e\bigl(x_n + g^o h\, t\bigr) = g^o\bigl(x_0 + h^{e+1} t\bigr) + q\,K(g, h),
\]
which makes explicit that the family of valid $(x_0, x_n)$ pairs forms an
arithmetic progression with step $(h^{e+1}, g^o h)$, motivating the modulus
$h^{e+1}$ in Step~4.

\medskip
\noindent\textbf{Step 4: Instantiate to concrete $(g, h)$ and solve for $x_0$.}
The function \texttt{path\_identity\_params($\sigma$, $g$, $h$)} substitutes the
concrete multiplier and base, computes $K$, and solves
\[
  x_0 \;\equiv\; -q\,K \cdot (g^o)^{-1} \pmod{h^{e+1}}
\]
for the canonical starting residue $x_0 \in [0, h^{e+1})$ whose trajectory follows
$\sigma$ for exactly $o + e$ steps.
The modulus is $h^{e+1}$ rather than $h^e$ because the word ends in $\OO$: after
all $e$ halvings the final odd step imposes one additional divisibility condition
on $x_0$.
For $h = 2$ (the standard Collatz system) this residue is unique; for $h \geq 3$
there may be multiple valid starting residues modulo $h^{e+1}$, and the formula
gives the canonical (smallest non-negative) representative.

\medskip
\noindent\textbf{Step 5: Read off the approximation data.}
The output tuple $(o,\, e,\, e/o,\, x_0)$ encodes:
\begin{itemize}
  \item $e/o$ --- a rational approximant to $\theta$, with $e/(o+e)$ and $o/(o+e)$
    the even and odd step densities;
  \item $x_0$ --- the canonical starting residue $(\bmod\; h^{e+1})$ whose $gx+q$
    trajectory realises the mediant word $\sigma$ for exactly $o + e$ steps.
\end{itemize}
Setting $t = 0$ gives $x_n = (g^o x_0 + qK) / h^e$, which is also determined by
$x_0$.
As $n \to \infty$ the ratio $e_{2n+1}/o_{2n+1}$ converges to $\theta$, with
convergence rate governed by the continued-fraction expansion of $\theta$.

\medskip
\noindent\textbf{Remark.}
Evaluating the Collatz iteration from $x_0$ is \emph{not} an efficient way to
generate the parity sequence: it requires arithmetic with integers of size
$O(h^e)$, which grows exponentially in the word length.
The parity sequence itself is produced directly and efficiently by the mediant
tree, with no large integer arithmetic.
The value of $x_0$ lies in establishing that a concrete Collatz trajectory
realising the sequence exists --- not in using it to regenerate the sequence.

\medskip
\noindent\textbf{Remark (residue constraint).}
To the extent that a $3x+1$ trajectory matches a factor of the Sturmian word for
$\theta = \log_2 3$, every odd element of the trajectory is congruent to $1$ or
$3$ modulo~$8$.
This is a consequence of the Sturmian alphabet being $\{\EE\OO, \EE\EE\OO\}$:
an odd step from a residue $r \equiv 1 \pmod{8}$ produces $v_2(3r+1) = 2$ even
steps (an $\EE\EE\OO$ block), while $r \equiv 3 \pmod{8}$ produces
$v_2(3r+1) = 1$ even step (an $\EE\OO$ block); residue $5 \pmod{8}$ forces the
parity sequence $\OO\EE\EE\EE\cdots$, and residue $7 \pmod{8}$ forces the parity
sequence $\OO\EE\OO\EE\OO$; neither pattern occurs in the Sturmian word for
$\log_2 3$, so neither residue can appear in a trajectory matching it.
A full account is given in~\cite{paper33}.

%------------------------------------------------------------
\section{The Greedy Parity Sequence Algorithm}
\label{sec:greedy}
%------------------------------------------------------------

For any generalised Collatz system $(g, h)$ with $g, h \geq 2$ and $\gcd(g,h)=1$,
there is a simple greedy algorithm that generates the idealised parity sequence
using only integer arithmetic --- no floating-point evaluation of $\log_h g$ is
required.
Let $\mathit{short} = \lfloor\log_h g\rfloor$ and
$\mathit{long} = \mathit{short} + 1$ (the two possible $\EE$-run lengths, since
$\log_h g \in (\mathit{short}, \mathit{long})$).

\begin{verbatim}
def collatz_log_h_of_g(g=3, h=2):
    e=o=0
    short=math.floor(math.log(g, h))
    long=short+1
    while True:
        if g**(o+1) < h**(e+short):
            yield from 'E' * short
            yield 'O'
            e+=short
            o+=1
        else:
            yield from 'E' * long
            yield 'O'
            e+=long
            o+=1
\end{verbatim}

The condition $g^{o+1} < h^{e+\mathit{short}}$ is a pure integer comparison
equivalent to $(o+1)\log_h g < e + \mathit{short}$, i.e.\ appending a
\emph{short} block $\EE^{\mathit{short}}\OO$ is sufficient to keep the ratio
$e/o$ above $\log_h g$; otherwise a \emph{long} block $\EE^{\mathit{long}}\OO$
is needed.
This is exactly the Beatty/Sturmian greedy rule: at each $\OO$ step the algorithm
selects the shorter run when the fractional part of $(o+1)\log_h g$ is already
covered by the accumulated $e$, and the longer run otherwise.

The resulting sequence is the Sturmian word for slope
$\alpha = 1/(1 + \log_h g)$, as verified by the identity
$s_n = \lfloor(n+1)\alpha\rfloor - \lfloor n\alpha\rfloor$ for all $n$.

For the standard case $g = 3$, $h = 2$ the first 60 symbols are:
\[
  \mathtt{EEOEEOEOEEOEOEEOEEOEOEEOEOEEOEEOEOEEOEOEEOEOEEOEEOEOEEOEOEEO},
\]
with $\OO$-density converging to $1/(1+\log_2 3) = \log_6 2 \approx 0.387$ ---
the \emph{ideally balanced} parity pattern, descending as quickly as theory
permits.
For $g = 5$, $h = 2$ ($\mathit{short} = 2$, $\mathit{long} = 3$, since
$\log_2 5 \approx 2.322$) the first 60 symbols are:
\[
  \mathtt{EEEOEEOEEOEEEOEEOEEOEEEOEEOEEOEEEOEEOEEOEEEOEEOEEOEEEOEEOEEO},
\]
with $\OO$-density converging to $1/(1+\log_2 5) \approx 0.301$.
The $5x+1$ system has known non-trivial cycles; the mediant word describes the
idealised parity pattern of a trajectory that descends at the asymptotically
optimal rate for that system.

%------------------------------------------------------------
\section{General Collatz Sequences are not Sturmian}
\label{sec:not-sturmian}
%------------------------------------------------------------

The Sturmian word for slope $\alpha = 1/(1+\theta)$ consists of exactly two block
types: $\EE^{\lfloor\theta\rfloor}\OO$ and $\EE^{\lceil\theta\rceil}\OO$,
interleaved in the Sturmian pattern.
These are precisely the leaf words of the mediant tree for $\theta$.
A Sturmian word can never contain a block of any other length.

A general Collatz trajectory is far less constrained: an actual orbit can produce
$\EE$-runs of any length, interleaved freely.
The mediant word (equivalently, the Sturmian word) is the \emph{ideally balanced}
special case.

Moreover, even among trajectories that approximate $\log_h g$ closely, the
Sturmian alphabet $\{\EE^{\lfloor\theta\rfloor}\OO,\, \EE^{\lceil\theta\rceil}\OO\}$
is not the only primitive alphabet that can arise.
For example, the known non-trivial cycles of the $5x+1$ system have alphabets
$\{\EE\OO, \EE\EE\EE\EE\OO\}$, $\{\EE\OO, \EE\EE\EE\EE\EE\OO\}$, and
$\{\EE\OO, \EE\EE\EE\OO\}$ --- each containing the block $\EE\OO$ of length $2$,
which does not appear in the Sturmian alphabet $\{\EE\EE\OO, \EE\EE\EE\OO\}$ at
all.
The mediant tree generates parity sequences over the Sturmian alphabet only; it
produces the optimal approximation to $\theta$ but does not exhaust all primitive
alphabets whose $e/o$ ratio approximates $\theta$.

%------------------------------------------------------------
\section{Relationship to Eliahou's Work}
\label{sec:eliahou}
%------------------------------------------------------------

Eliahou~\cite{eliahou1993} studies rational approximation of $\log_2 3$ in the
context of bounding nontrivial cycle lengths in the Collatz problem.
His approach identifies best upper rational approximants via the continued-fraction
expansion of $\log_2 3$, yielding explicit lower bounds on Collatz cycle lengths.

Our work differs in two respects.
First, our algorithm constructs approximants directly from bounded pairs via the
optimal weight formula (Theorem~8.1 of~\cite{paper76}), with no continued-fraction
computation; the connection to continued fractions is a consequence of the Sturmian
identity.
Second, and more fundamentally, Eliahou's framework gives no means of constructing
a parity sequence realising a given approximant.
The mediant tree simultaneously produces the rational approximant and the mediant
word that realises it, as a byproduct of the same construction.

Where Eliahou's work illuminates the relationship between Collatz trajectories and
good rational approximations of $\log_2 3$ specifically, the present work
generalises in both directions: the mediant tree constructs good rational
approximants to \emph{any} real $\theta$ (see~\cite{paper76}), and the pipeline
of Section~\ref{sec:pipeline} produces a concrete $gx+q$ trajectory (for any
$g, h \geq 2$ with $\gcd(g,h)=1$ and $q \neq 0$) realising each approximant.
The result is that every sufficiently good rational approximation to any real
number is realisable by a constructible generalised Collatz trajectory.

%------------------------------------------------------------
\begin{thebibliography}{9}

\bibitem{paper76}
  J.~Seymour,
  \textit{Rational Approximations via the Mediant Tree},
  \texttt{https://github.com/wildducktheories/collatz}, Paper~76 (2026).

\bibitem{eliahou1993}
  S.~Eliahou,
  \textit{The $3x+1$ problem: new lower bounds on nontrivial cycle lengths},
  Discrete Mathematics \textbf{118} (1993), 45--56.

\bibitem{paper71}
  J.~Seymour,
  \textit{The Structural Polynomial of a Parity Sequence},
  \texttt{https://github.com/wildducktheories/collatz}, Paper~71 (2026).

\bibitem{paper33}
  J.~Seymour,
  \textit{A Regular Expression Language for the Collatz Graph},
  \texttt{https://github.com/wildducktheories/collatz}, Paper~33 (2026).

\end{thebibliography}

\end{document}
